
\chapter{معرفی محصول و مسئله}

\section{InsightInk چیست؟}

\eng{InsightInk} یک سامانه محتوایی است که ایده‌های مفید استخراج‌شده از کتاب، مقاله، فیلم/مستند و پادکست را در قالب یک Feed در اختیار کاربر قرار می‌دهد. ایده اصلی محصول این است که «اسکرول کردن» به‌جای مصرف محتوای صرفاً سرگرم‌کننده، به مواجهه با تجربه‌ها و ایده‌های قابل‌فکر و قابل‌عمل تبدیل شود.

\begin{sourcebox}
در PDS، محصول به‌عنوان فضایی برای مرور ایده‌های حاصل از تجربه خواندن کتاب و مقاله، دیدن فیلم و شنیدن پادکست معرفی شده است. کاربر می‌تواند از این محتوا الهام بگیرد، آن را ذخیره یا به اشتراک بگذارد و با سازندگان محتوا ارتباط بگیرد.
\end{sourcebox}

\section{هدف این پروژه دانشگاهی}

هدف ما ساخت یک استارتاپ production در مقیاس میلیون‌ها کاربر نیست. هدف عبارت است از:

\begin{itemize}
  \item ساخت یک پروژه Full-Stack واقعی و قابل اجرا؛
  \item یادگیری عمیق \eng{Python} و \eng{Django}؛
  \item یادگیری Frontend از مبانی تا \eng{React + TypeScript}؛
  \item طراحی درست دیتابیس رابطه‌ای؛
  \item ساخت API، احراز هویت، Permission و تست؛
  \item پیاده‌سازی سیستم پیشنهاددهی الگوریتمی قابل توضیح؛
  \item پشتیبانی هم‌زمان از فارسی و انگلیسی؛
  \item داشتن مستندات کامل برای ادامه کار بدون وابستگی به AI.
\end{itemize}

\begin{decisionbox}
معماری باید برای یک نفر قابل نگهداری باشد. بنابراین \eng{microservice}، \eng{Kubernetes}، \eng{event bus}، \eng{vector database} و زیرساخت‌های سنگین وارد Baseline نمی‌شوند.
\end{decisionbox}

\section{Actorهای اصلی}

\subsection{Reader / User}

کاربر معمولی که می‌تواند:
\begin{itemize}
  \item ایده‌ها را بخواند؛
  \item Spark بزند؛
  \item ذخیره کند؛
  \item نظر بدهد؛
  \item Creator را Follow کند؛
  \item ایده‌ها را جست‌وجو کند؛
  \item کتابخانه شخصی داشته باشد؛
  \item محتوای نامناسب را Report کند؛
  \item در فاز پیام‌رسانی با Creator تماس بگیرد.
\end{itemize}

\subsection{Creator}

Creator یک User جداگانه از نظر هویتی نیست؛ یک User است که \eng{CreatorProfile} و Permissionهای مربوط را دارد. او می‌تواند پروفایل سازنده داشته باشد، محتوای خودش را منتشر/مدیریت کند و با مخاطب تعامل داشته باشد.

\subsection{Moderator}

Moderator ناظر سیستم است و دسترسی‌های محدود اما قوی برای Trust \& Safety دارد:
\begin{itemize}
  \item مشاهده Reportها؛
  \item بررسی محتوای گزارش‌شده؛
  \item مخفی/حذف کردن Comment یا Content طبق سیاست؛
  \item ثبت اقدام Moderation؛
  \item اعمال محدودیت‌های مشخص روی User در صورت نیاز.
\end{itemize}

Moderator مدیر کل نیست و اجازه ایجاد Moderator یا Admin ندارد.

\subsection{Admin}

Admin مدیر کل سیستم است و بالاترین دسترسی را دارد. فقط Admin می‌تواند Moderatorها را تعیین یا حذف کند.

\section{زبان‌های محصول}

\begin{decisionbox}
زبان اصلی محصول فارسی و زبان دوم انگلیسی است. از روز اول ساختار Frontend و مدل داده باید امکان \eng{RTL/LTR} و محتوای دو زبانه را در نظر بگیرد.
\end{decisionbox}

این یعنی:
\begin{itemize}
  \item UI فارسی راست‌به‌چپ؛
  \item UI انگلیسی چپ‌به‌راست؛
  \item زبان UI از زبان Content جداست؛
  \item کاربر می‌تواند UI فارسی داشته باشد و محتوای انگلیسی ببیند.
\end{itemize}

\section{تعریف AI در این پروژه}

\begin{warningbox}
در منابع اولیه عباراتی مانند \eng{AI-Assisted Tool} و \eng{AI-Powered Idea Discovery} آمده است. در Scope نهایی این پروژه، هیچ LLM یا RAG یا Agent وجود ندارد.
\end{warningbox}

سیستم پیشنهاددهی ما کاملاً الگوریتمی است و با داده‌هایی مانند Interest، Save، Spark، Follow، تازگی محتوا و Popularity کار می‌کند.

مواردی که صریحاً خارج از پروژه هستند:
\begin{itemize}
  \item \eng{OpenAI/Claude/Gemini API}
  \item \eng{LLM}
  \item \eng{RAG}
  \item \eng{Embeddings}
  \item \eng{Vector DB}
  \item \eng{Prompt Engineering}
  \item \eng{AI Agent}
\end{itemize}

\section{معیار موفقیت پروژه}

در پایان، پروژه باید حداقل بتواند Flow زیر را واقعاً اجرا کند:

\begin{enumerate}
  \item کاربر Account بسازد و Login کند؛
  \item زبان و علاقه‌مندی‌ها را انتخاب کند؛
  \item Feed واقعی از PostgreSQL ببیند؛
  \item Idea را باز کند؛
  \item Source آن را ببیند؛
  \item Spark/Save/Comment کند؛
  \item Saved Ideaها را در Library ببیند؛
  \item Creator را ببیند و Follow کند؛
  \item Search/Explore داشته باشد؛
  \item Recommendation الگوریتمی دریافت کند؛
  \item UI فارسی و انگلیسی درست داشته باشد.
\end{enumerate}

Realtime، Messaging، Notification، TTS و Membership/eBook بعد از این Core Flow اضافه می‌شوند.
